% unit: noether.p13.sec.004
\section*{\S{}4. Réciproque dans le cas du groupe infini}

% unit: noether.p13.par.061
Il faut d'abord montrer que l'hypothèse de linéarité de \(\D x\) et de \(\D u\) ne constitue pas une restriction. Ici cela résulte, même sans la réciproque, du fait que \(\G_{\infty\rho}\) dépend formellement de \(\rho\), et seulement de \(\rho\), fonctions arbitraires. En effet, dans le cas non linéaire, lors de la composition des transformations, les termes de plus bas ordre s'additionnent et le nombre des fonctions arbitraires augmenterait. Supposons par exemple que
% unit: noether.p13.eq.018a
\[
\begin{aligned}
y&=A\!\left(x,u,\frac{\p u}{\p x},\ldots;p\right)\\
&=x+\sum a(x,u,\ldots)p^\nu
+b(x,u,\ldots)p^{\nu-1}\frac{\p p}{\p x}
+c\,p^{\nu-2}\left(\frac{\p p}{\p x}\right)^2+\cdots
+d\left(\frac{\p p}{\p x}\right)^\nu+\cdots,
\end{aligned}
\]
où
% unit: noether.p13.eq.018b
\[
p^\nu=(p^{(1)})^{\nu_1}\cdots(p^{(\rho)})^{\nu_\rho},
\]
et, de manière correspondante,
% unit: noether.p13.eq.018c
\[
v=B\!\left(x,u,\frac{\p u}{\p x},\ldots;p\right).
\]
Alors, par composition avec
% unit: noether.p13.eq.018d
\[
z=A\!\left(y,v,\frac{\p v}{\p y},\ldots;q\right),
\]
les termes de plus bas ordre donnent
% unit: noether.p13.eq.018e
\[
z=x+\sum a(p^\nu+q^\nu)+
b\left\{p^{\nu-1}\frac{\p p}{\p x}+q^{\nu-1}\frac{\p q}{\p x}\right\}
+c\left\{p^{\nu-2}\left(\frac{\p p}{\p x}\right)^2
+q^{\nu-2}\left(\frac{\p q}{\p x}\right)^2\right\}+\cdots.
\]
Si, ici, un coefficient différent de \(a\) et de \(b\) est non nul, de sorte qu'un terme
% unit: noether.p13.eq.018f
\[
p^{\nu-\sigma}\left(\frac{\p p}{\p x}\right)^\sigma
+q^{\nu-\sigma}\left(\frac{\p q}{\p x}\right)^\sigma
\]
apparaisse effectivement pour \(\sigma>1\), celui-ci ne peut pas s'écrire comme quotient différentiel d'une seule fonction, ni comme produit de puissances d'un tel quotient. Le nombre des fonctions arbitraires a donc augmenté, contrairement à l'hypothèse. Si tous les coefficients autres que \(a\) et \(b\) s'annulent, alors, suivant les valeurs des exposants \(\nu_1,\ldots,\nu_\rho\), le second terme est le quotient différentiel du premier (comme par exemple toujours pour une \(\G_{\infty1}\)), de sorte que la linéarité se présente effectivement; ou bien le nombre des fonctions arbitraires augmente également dans ce cas. Les transformations infinitésimales satisfont donc, en raison de la linéarité en \(p(x)\), à un système d'équations différentielles partielles linéaires; et, puisque la propriété de groupe est vérifiée, elles forment un \og groupe infini de transformations infinitésimales \fg{} au sens de Lie (\emph{Grundlagen}, \S{}10).

% unit: noether.p13.par.062
La réciproque s'obtient maintenant de manière analogue au cas du groupe fini. L'existence des dépendances (16), après multiplication par \(p^{(\lambda)}(x)\) et addition, conduit, au moyen de la transformation identique (14), à
% unit: noether.p13.eq.019
\[
\sum_i\psi_i\bd u_i=\Div\Gamma.
\]
De là on obtient, comme au \S{}3, la détermination de \(\D x\) et de \(\D u\), ainsi que l'invariance de \(I\) relativement à ces transformations infinitésimales, qui dépendent effectivement linéairement de \(\rho\) fonctions arbitraires et de leurs dérivées jusqu'à l'ordre \(\sigma\). Que ces transformations infinitésimales, lorsqu'elles ne contiennent pas les dérivées \(\p u/\p x,\ldots\), forment sûrement un groupe, résulte comme au \S{}3: sinon, par composition, apparaîtraient plus de fonctions arbitraires, alors que, par hypothèse, il ne doit y avoir que \(\rho\) dépendances (16). Elles forment donc un \og groupe infini de transformations infinitésimales \fg{}. Or un tel groupe consiste (\emph{Grundlagen}, théorème VII, p. 391) en les transformations infinitésimales les plus générales d'un certain \og groupe infini \(\G\) de transformations finies \fg{} ainsi défini au sens de Lie. Chaque transformation finie est alors engendrée à partir de transformations infinitésimales (\emph{Grundlagen}, \S{}7),\footnote{Il en résulte en particulier que le groupe \(\G\) engendré à partir des transformations infinitésimales \(\D x,\D u\) d'une \(\G_{\infty\rho}\) ramène de nouveau à \(\G_{\infty\rho}\). Car \(\G_{\infty\rho}\) ne contient pas de transformations infinitésimales dépendant de fonctions arbitraires et différentes de \(\D x,\D u\), et ne peut pas non plus contenir de transformations indépendantes dépendant de paramètres, puisqu'autrement elle serait un groupe mixte. Or, d'après ce qui précède, les transformations finies sont déterminées par les infinitésimales.} et provient donc de l'intégration du système simultané:
% unit: noether.p13.eq.020
\[
\frac{\dd x_i}{\dd t}=\D x_i,\qquad
\frac{\dd u_i}{\dd t}=\D u_i,\qquad
\left\{\begin{array}{l}
x_i=y_i,\\
u_i=v
\end{array}\right\}
\quad\text{pour }t=0.
\]
Il peut toutefois être nécessaire de laisser les \(p(x)\) arbitraires dépendre encore de \(t\). Ainsi \(\G\) dépend effectivement de \(\rho\) fonctions arbitraires. En particulier, s'il suffit de prendre \(p(x)\) indépendante de \(t\), cette dépendance est analytique dans les fonctions arbitraires \(q(x)=t\cdot p(x)\).\footnote{La question de savoir si ce dernier cas se produit peut-être toujours a été posée par Lie dans une autre formulation (\emph{Grundlagen}, \S{}7 et \S{}13, fin).} Si des dérivées \(\p u/\p x,\ldots\) interviennent, il peut être nécessaire d'adjoindre encore des transformations infinitésimales \(\bd u=0,\ \Div(f\cdot\D x)=0\) avant de tirer les mêmes conclusions.

% unit: noether.p13.par.063
À la suite d'un exemple de Lie (\emph{Grundlagen}, \S{}7), indiquons encore un cas assez général où l'on peut aller jusqu'à des formules explicites, qui montrent en même temps que les dérivées des fonctions arbitraires apparaissent jusqu'à l'ordre \(\sigma\); ici la réciproque est donc complète. Il s'agit des groupes de transformations infinitésimales auxquels correspond le groupe de toutes les transformations des \(x\), et des transformations des \(u\) qui en sont \og induites \fg{}; c'est-à-dire des transformations des \(u\) dans lesquelles \(\D u\), et par conséquent \(u\), ne dépend que des fonctions arbitraires qui apparaissent dans \(\D x\). On suppose encore que les dérivées \(\p u/\p x,\ldots\) n'apparaissent pas dans \(\D u\). On a donc:
% unit: noether.p13.eq.021
\[
\D x_i=p^{(i)}(x),\qquad
\D u_i=
\sum_{\lambda=1}^{n}\left\{
a^{(\lambda)}(x,u)p^{(\lambda)}
+b^{(\lambda)}\frac{\p p^{(\lambda)}}{\p x}
+\cdots+
c^{(\lambda)}\frac{\p^\sigma p^{(\lambda)}}{\p x^\sigma}
\right\}.
\]
Comme la transformation infinitésimale \(\D x=p(x)\) engendre toute transformation \(x=y+g(y)\) avec \(g(y)\) arbitraire, on peut en particulier déterminer \(p(x)\) en fonction de \(t\) de manière que soit engendré le groupe à un paramètre:
% unit: noether.p13.eq.022
\begin{equation}
x_i=y_i+t\,g_i(y),
\tag{18}
\end{equation}
qui, pour \(t=0\), devient l'identité et, pour \(t=1\), devient la transformation cherchée \(x=y+g(y)\). En différentiant (18), on obtient en effet:
% unit: noether.p13.eq.023
\begin{equation}
\frac{\dd x_i}{\dd t}=g_i(y)=p^{(i)}(x,t),
\tag{19}
\end{equation}
où \(p(x,t)\) est déterminé à partir de \(g(y)\) par inversion de (18); inversement, (18) résulte de (19) au moyen de la condition auxiliaire \(x_i=y_i\) pour \(t=0\), qui fixe l'intégrale de manière univoque. Au moyen de (18), les \(x\) qui figurent dans \(\D u\) peuvent être remplacés par les constantes d'intégration \(y\) et par \(t\); les \(g(y)\) apparaissent alors exactement jusqu'à leur \(\sigma\)-ième dérivée, car dans
% unit: noether.p13.eq.024
\[
\frac{\p p}{\p x}=\sum\frac{\p g}{\p y_\varkappa}\frac{\p y_\varkappa}{\p x}
\]
les \(\p y/\p x\) sont exprimées par \(\p x/\p y\), et, en général, \(\p^\sigma p/\p x^\sigma\) est remplacé par sa valeur dans
% unit: noether.p13.eq.025
\[
\frac{\p g}{\p y},\ldots,\frac{\p x}{\p y},\ldots,\frac{\p^\sigma x}{\p y^\sigma}.
\]
Pour déterminer les \(u\), on obtient donc le système d'équations
% unit: noether.p13.eq.026
\[
\frac{\dd u_i}{\dd t}
=
F_i\!\left(g(y),\frac{\p g}{\p y},\ldots,\frac{\p^\sigma g}{\p y^\sigma},u,t\right),
\qquad (u_i=v_i\ \text{pour }t=0),
\]
où seules \(t\) et \(u\) sont variables, tandis que \(g(y),\ldots\) appartiennent au domaine des coefficients; l'intégration donne donc
% unit: noether.p13.eq.027
\[
u_i=v_i+
B_i\!\left(v,g(y),\frac{\p g}{\p y},\ldots,\frac{\p^\sigma g}{\p y^\sigma},t\right)_{t=1}.
\]
On obtient ainsi des transformations qui dépendent exactement de \(\sigma\) dérivées des fonctions arbitraires. L'identité y est contenue pour \(g(y)=0\), d'après (18); et la propriété de groupe résulte du fait que le procédé indiqué fournit chaque transformation \(x=y+g(y)\), ce par quoi la transformation induite des \(u\) est fixée univoquement, de sorte que le groupe \(\G\) est épuisé.

% unit: noether.p13.par.064
Il résulte encore incidemment de la réciproque que ce n'est pas une restriction de supposer les fonctions arbitraires dépendantes seulement des \(x\), et non de \(u,\p u/\p x,\ldots\). Dans ce dernier cas, dans la transformation identique (14), donc aussi dans (15), apparaîtraient, en plus des \(p^{(\lambda)}\), les dérivées \(\p p^{(\lambda)}/\p u,\ \p p^{(\lambda)}/\p(\p u/\p x),\ldots\). Si l'on prend alors successivement les \(p^{(\lambda)}\) comme étant de degré zéro, premier, etc. en \(u,\p u/\p x,\ldots\), avec des fonctions arbitraires de \(x\) comme coefficients, des dépendances (16) apparaissent de nouveau, mais en plus grand nombre; par la réciproque précédente, elles ramènent toutefois au cas antérieur par regroupement avec des fonctions arbitraires qui ne dépendent que de \(x\). De même, on montre qu'au surgissement simultané de dépendances et de relations de divergence indépendantes de celles-ci correspondent des groupes mixtes.\footnote{Comme au \S{}3, il résulte encore ici de la réciproque qu'à côté de \(I\), toute intégrale \(I^*\) qui diffère de lui par une intégrale portant sur une divergence admet également un groupe infini avec les mêmes \(\bd u\); ici toutefois \(\D x\) et \(\D u\) contiendront en général des dérivées des \(u\). Einstein a introduit une telle intégrale \(I^*\) dans la relativité générale afin d'obtenir une formulation plus simple des théorèmes d'énergie; j'indique les transformations infinitésimales que cet \(I^*\) admet, en me conformant exactement à la notation de la seconde note de Klein. L'intégrale \(I=\int\!\cdots\!\int K\,\dd\omega=\int\!\cdots\!\int \mathfrak{R}\,\dd S\) admet le groupe de toutes les transformations des \(w\) et le groupe induit par là pour les \(g_{\mu\nu}\). À cela correspondent les dépendances ((30) chez Klein):
\[
\sum \mathfrak{R}_{\mu\nu}g_\tau^{\mu\nu}
+2\sum\frac{\p g^{\mu\nu}\mathfrak{R}_{\mu\tau}}{\p w^\sigma}=0.
\]
Or \(I^*=\int\!\cdots\!\int \mathfrak{R}^*\,\dd S\), où \(\mathfrak{R}^*=\mathfrak{R}+\Div\), et par conséquent \(\mathfrak{R}_{\mu\nu}^*=\mathfrak{R}_{\mu\nu}\), les \(\mathfrak{R}_{\mu\nu}^*,\mathfrak{R}_{\mu\nu}\) désignant respectivement les expressions de Lagrange. Les dépendances indiquées sont donc aussi des dépendances pour \(\mathfrak{R}_{\mu\nu}^*\); et, après multiplication par \(p^\tau\) et addition, en renversant les transformations de la différentiation d'un produit, on obtient
\[
\sum \mathfrak{R}_{\mu\nu}p^{\mu\nu}
+2\Div\!\left(\sum g^{\mu\sigma}\mathfrak{R}_{\mu\tau}p^\tau\right)=0,
\]
\[
\delta\mathfrak{R}^*+\Div\!\left(\sum 2g^{\mu\sigma}\mathfrak{R}_{\mu\tau}p^\tau
-\frac{\p\mathfrak{R}^*}{\p g^{\mu\nu}_{\sigma}}p^{\mu\nu}\right)=0.
\]
La comparaison avec l'équation différentielle de Lie
\(\delta\mathfrak{R}^*+\Div(\mathfrak{R}^*\D w)=0\) donne
\[
\D w^\sigma=\frac1{\mathfrak{R}^*}\cdot\left(\sum 2g^{\mu\sigma}\mathfrak{R}_{\mu\tau}p^\tau
-\frac{\p\mathfrak{R}^*}{\p g^{\mu\nu}_{\sigma}}p^{\mu\nu}\right),
\qquad
\D g^{\mu\nu}=p^{\mu\nu}+\sum g^{\mu\nu}_{\sigma}\D w^\sigma,
\]
comme transformations infinitésimales admises par \(I^*\). Ces transformations infinitésimales dépendent donc des premières et secondes dérivées des \(g^{\mu\nu}\), et contiennent les \(p\) arbitraires jusqu'à la première dérivée.}

% unit: noether.p13.sec.005
\section*{\S{}5. Invariance des composantes individuelles des relations}

% unit: noether.p13.par.065
Si l'on spécialise le groupe \(\G\) au cas le plus simple, ordinairement considéré, en n'admettant dans les transformations aucune dérivée des \(u\), et en faisant dépendre les variables indépendantes transformées seulement des \(x\), non des \(u\), alors on peut conclure à l'invariance des composantes individuelles dans les formules. D'abord, par des raisonnements connus, on obtient l'invariance de
% unit: noether.p13.eq.028
\[
\int\!\cdots\!\int\left(\sum\psi_i\delta u_i\right)\dd x,
\]
donc l'invariance relative de \(\sum\psi_i\delta u_i\),\footnote{C'est-à-dire que \(\sum\psi_i\delta u_i\) acquiert un facteur sous les transformations, ce que la théorie algébrique des invariants a toujours appelé invariance relative.} où \(\delta\) désigne une variation quelconque. En effet, d'une part,
% unit: noether.p13.eq.029
\[
\delta I
=\int\!\cdots\!\int \delta f\!\left(x,u,\frac{\p u}{\p x},\ldots\right)\dd x
=\int\!\cdots\!\int \delta f\!\left(y,v,\frac{\p v}{\p y},\ldots\right)\dd y.
\]
D'autre part, pour \(\delta u,\delta(\p u/\p x),\ldots\) s'annulant au bord, auxquels, en raison de la transformation linéaire homogène des \(\delta u,\delta(\p u/\p x),\ldots\), correspondent des \(\delta v,\delta(\p v/\p y),\ldots\) s'annulant également au bord, on a
% unit: noether.p13.eq.030
\[
\int\!\cdots\!\int \delta f\!\left(x,u,\frac{\p u}{\p x},\ldots\right)\dd x
=
\int\!\cdots\!\int\left(\sum\psi_i(u,\ldots)\delta u_i\right)\dd x,
\]
% unit: noether.p13.eq.031
\[
\int\!\cdots\!\int \delta f\!\left(y,v,\frac{\p v}{\p y},\ldots\right)\dd y
=
\int\!\cdots\!\int\left(\sum\psi_i(v,\ldots)\delta v_i\right)\dd y.
\]
Par suite, pour \(\delta u,\delta(\p u/\p x),\ldots\) s'annulant au bord,
% unit: noether.p13.eq.032
\[
\int\!\cdots\!\int\left(\sum\psi_i(u,\ldots)\delta u_i\right)\dd x
=
\int\!\cdots\!\int\left(\sum\psi_i(v,\ldots)\delta v_i\right)\dd y
\]
% unit: noether.p13.eq.033
\[
=
\int\!\cdots\!\int\left(\sum\psi_i(v,\ldots)\delta v_i\right)
\left|\frac{\p y_i}{\p x_\varkappa}\right|\dd x.
\]
Si l'on exprime dans la troisième intégrale \(y,v,\delta v\) par \(x,u,\delta u\), et qu'on l'égale à la première, on obtient donc une relation
% unit: noether.p13.eq.034
\[
\int\!\cdots\!\int\left(\sum\chi_i(u,\ldots)\delta u_i\right)\dd x=0
\]
pour \(\delta u\) s'annulant au bord, mais par ailleurs arbitraire. Il en résulte, comme on sait, l'annulation de l'intégrande pour \(\delta u\) arbitraire; la relation
% unit: noether.p13.eq.035
\[
\sum\psi_i(u,\ldots)\delta u_i
=
\left|\frac{\p y_i}{\p x_\varkappa}\right|
\left(\sum\psi_i(v,\ldots)\delta v_i\right)
\]
vaut donc identiquement en \(\delta u\). Elle exprime l'invariance relative de \(\sum\psi_i\delta u_i\), et par conséquent l'invariance de
% unit: noether.p13.eq.036
\[
\int\!\cdots\!\int\left(\sum\psi_i\delta u_i\right)\dd x.
\]
\footnote{Ces raisonnements échouent lorsque \(y\) dépend aussi des \(u\), car alors \(\delta f(y,v,\p v/\p y,\ldots)\) contient aussi des termes \(\sum(\p f/\p y)\delta y\), de sorte que la transformation de divergence ne conduit pas aux expressions de Lagrange. Ils échouent de même si l'on admet des dérivées des \(u\); alors les \(\delta v\) deviennent des combinaisons linéaires de \(\delta u,\delta(\p u/\p x),\ldots\), et ne conduisent donc qu'après une nouvelle transformation de divergence à une identité \(\int\!\cdots\!\int(\sum\chi_i(u,\ldots)\delta u_i)\dd x=0\), si bien qu'à droite les expressions de Lagrange n'apparaissent de nouveau pas. La question de savoir si l'invariance de \(\int\!\cdots\!\int(\sum\psi_i\delta u_i)\dd x\) permet déjà de conclure à l'existence de relations de divergence est, d'après la réciproque, équivalente à la question de savoir si elle entraîne l'invariance de \(I\) vis-à-vis d'un groupe qui ne conduit pas nécessairement aux mêmes \(\D u,\D x\), mais bien aux mêmes \(\bd u\). Dans le cas spécial de l'intégrale simple et de seules dérivées premières dans \(f\), on peut, pour le groupe fini, conclure de l'invariance des expressions de Lagrange à l'existence de premières intégrales (voir par exemple Engel, Gött. Nachr. 1916, p. 270).}

% unit: noether.p13.par.066
Pour appliquer ceci aux relations de divergence et aux dépendances dérivées, il faut d'abord montrer que le \(\bd u\) déduit de \(\D u,\D x\) satisfait effectivement aux lois de transformation de la variation \(\delta u\), pourvu seulement que les paramètres, respectivement les fonctions arbitraires, dans \(\bd v\) soient déterminés comme ils correspondent au groupe semblable des transformations infinitésimales en \(y,v\). Soit \(\mathfrak T_q\) la transformation qui fait passer \(x,u\) en \(y,v\), et soit \(\mathfrak T_p\) une transformation infinitésimale en \(x,u\). Alors la transformation semblable en \(y,v\) est donnée par \(\mathfrak T_r=\mathfrak T_q\mathfrak T_p\mathfrak T_q^{-1}\), les paramètres, respectivement les fonctions arbitraires, \(r\) étant donc déterminés par \(p\) et \(q\). En formules, cela s'exprime ainsi:
% unit: noether.p13.eq.037
\[
\mathfrak T_p:\ \xi=x+\D x(x,p),\qquad u^*=u+\D u(x,u,p);
\]
% unit: noether.p13.eq.038
\[
\mathfrak T_q:\ y=A(x,q),\qquad v=B(x,u,q);
\]
% unit: noether.p13.eq.039
\[
\mathfrak T_q\mathfrak T_p:\ \eta=A(x+\D x(x,p),q),\quad
v^*=B(x+\D x(p),\,u+\D u(p),q).
\]
Il en résulte \(\mathfrak T_r=\mathfrak T_q\mathfrak T_p\mathfrak T_q^{-1}\), donc
% unit: noether.p13.eq.040
\[
\eta=y+\D y(r),\qquad v^*=v+\D v(r),
\]
où, au moyen de l'inverse de \(\mathfrak T_q\), on considère les \(x\) comme fonctions des \(y\) et l'on ne retient que les termes infinitésimaux. On a ainsi l'identité
% unit: noether.p13.eq.041
\begin{equation}
\begin{gathered}
\eta=y+\D y(r)=y+\sum\frac{\p A(x,q)}{\p x}\D x(p),\\[0.55em]
v^*=v+\D v(r)=v+\sum\frac{\p B(x,u,q)}{\p x}\D x(p)
+\sum\frac{\p B(x,u,q)}{\p u}\D u(p).
\end{gathered}
\tag{20}
\end{equation}
Si l'on y remplace \(\xi=x+\D x\) par \(\xi-\D \xi\), de sorte que \(\xi\) revienne à \(x\), et que donc \(\D x\) s'annule, alors, d'après la première formule (20), \(\eta\) revient aussi à \(y=\eta-\D\eta\). Si, sous cette substitution, \(\D u(p)\) passe en \(\bd u(p)\), alors \(\D v(r)\) passe aussi en \(\bd v(r)\), et la seconde formule (20) donne
% unit: noether.p13.eq.043
\[
v+\bd v(y,v,\ldots,r)=v+\sum\frac{\p B(x,u,q)}{\p u}\,\bd u(p),
\]
% unit: noether.p13.eq.044
\[
\bd v(y,v,\ldots,r)=\sum\frac{\p B}{\p u_\varkappa}\,\bd u_\varkappa(x,u,p).
\]
Les formules de transformation pour les variations sont donc effectivement satisfaites, pourvu seulement que \(\bd v\) soit pris comme dépendant des paramètres, respectivement des fonctions arbitraires, \(r\).\footnote{On voit de nouveau qu'il faut supposer \(y\) indépendant de \(u\), etc., pour que les conclusions soient valables. Comme exemple, on peut citer les \(\delta g^{\mu\nu}\) et \(\delta q_\sigma\) indiqués par Klein, qui satisfont aux transformations pour variations dès que les \(p\) sont soumis à une transformation vectorielle.}

% unit: noether.p13.par.067
Il s'ensuit en particulier l'invariance relative de \(\sum\psi_i\bd u_i\); donc aussi, d'après (12), puisque les relations de divergence sont également satisfaites en \(y,v\), l'invariance relative de \(\Div\bar B\); et encore, d'après (14) et (13), l'invariance relative de \(\Div\Gamma\) et des membres de gauche des dépendances regroupés avec les \(p^{(\lambda)}\), où, dans les formules transformées, les \(p(x)\) arbitraires (respectivement les paramètres) doivent partout être remplacés par les \(r\). Il en résulte encore l'invariance relative de \(\Div(B-\Gamma)\), c'est-à-dire d'une divergence d'un système de fonctions \(B-\Gamma\) non identiquement nul, mais dont la divergence s'annule identiquement.

% unit: noether.p13.par.068
De l'invariance relative de \(\Div B\), dans le cas unidimensionnel et pour le groupe fini, on peut encore tirer une conclusion sur l'invariance des premières intégrales. La transformation des paramètres correspondant à la transformation infinitésimale devient, d'après (20), linéaire et homogène; et, parce que toutes les transformations sont inversibles, les \(\eps\) sont aussi linéaires et homogènes dans les paramètres transformés \(\eps^*\). Cette inversibilité demeure certainement conservée lorsqu'on pose \(\psi=0\), puisque aucune dérivée de \(u\) n'intervient dans (20). En égalant les coefficients des \(\eps^*\) dans
% unit: noether.p13.eq.045
\[
\Div B(x,u,\ldots,\eps)=\frac{\dd y}{\dd x}\,\Div B(y,v,\ldots,\eps^*)
\]
les quantités \(\frac{\dd}{\dd y}B^{(\lambda)}(y,v,\ldots)\) deviennent aussi des fonctions linéaires homogènes des \(\frac{\dd}{\dd x}B^{(\lambda)}(x,u,\ldots)\). Ainsi, de
% unit: noether.p13.eq.046
\[
\frac{\dd}{\dd x}B^{(\lambda)}(x,u,\ldots)=0
\quad\text{ou}\quad B^{(\lambda)}(x,u)=\mathrm{const.}
\]
il résulte
% unit: noether.p13.eq.047
\[
\frac{\dd}{\dd y}B^{(\lambda)}(y,v,\ldots)=0
\quad\text{ou}\quad B^{(\lambda)}(y,v)=\mathrm{const.}
\]
Les \(\rho\) premières intégrales qui correspondent à une \(\G_\rho\) admettent donc chacune le groupe, de sorte que l'intégration ultérieure s'en trouve également simplifiée. L'exemple le plus simple en est le cas où \(f\) ne dépend pas de \(x\), ou d'un \(u\), ce qui correspond aux transformations infinitésimales \(\D x=\eps,\D u=0\), respectivement \(\D x=0,\D u=\eps\). Alors \(\bd u=-\eps\,\dd u/\dd x\), respectivement \(\eps\); et, comme \(B\) se déduit de \(f\) et de \(\bd u\) par différentiation et combinaisons rationnelles, il ne dépend donc pas non plus de \(x\), respectivement de \(u\), et admet les groupes correspondants.\footnote{Dans les cas où l'existence de premières intégrales résulte déjà de l'invariance de \(\int(\sum\psi_i\delta u_i)\dd x\), celles-ci n'admettent pas le groupe complet \(\G_\rho\). Par exemple, \(\int(u''\delta u)\dd x\) admet la transformation infinitésimale \(\D x=\eps_2,\D u=\eps_1+x\eps_3\), tandis que la première intégrale \(u-u'x=\mathrm{const.}\), correspondant à \(\D x=0,\D u=x\eps_3\), n'admet pas les deux autres transformations infinitésimales, puisqu'elle contient explicitement à la fois \(u\) et \(x\). À cette première intégrale correspondent précisément des transformations infinitésimales pour \(f\) qui contiennent des dérivées. On voit donc que l'invariance de \(\int\!\cdots\!\int(\sum\psi_i\delta u_i)\dd x\) réalise en tout cas moins que l'invariance de \(I\), remarque à faire à propos de la question soulevée dans une note précédente.}

% unit: noether.p13.sec.006
\section*{\S{}6. Une assertion de Hilbert}

% unit: noether.p13.par.069
De ce qui précède résulte enfin la preuve d'une assertion de Hilbert concernant le lien entre l'échec des théorèmes d'énergie propres et la \og relativité générale \fg{} (première note de Klein, Göttinger Nachr. 1917, Réponse, premier paragraphe), sous une forme généralisée du point de vue de la théorie des groupes.

% unit: noether.p13.par.070
Supposons que l'intégrale \(I\) admette une \(\G_{\infty\rho}\), et que \(\G_\sigma\) soit un groupe fini quelconque provenant d'une spécialisation des fonctions arbitraires, donc un sous-groupe de \(\G_{\infty\rho}\). Au groupe infini \(\G_{\infty\rho}\) correspondent alors des dépendances (16), au groupe fini \(\G_\sigma\) des relations de divergence (13); et, réciproquement, de l'existence de relations de divergence quelconques résulte l'invariance de \(I\) par rapport à un groupe fini, qui est identique à \(\G_\sigma\) si et seulement si les \(\bd u\) sont des combinaisons linéaires de celles qui proviennent de \(\G_\sigma\). L'invariance par rapport à \(\G_\sigma\) ne peut donc conduire à aucune relation de divergence différente de (13). Mais comme de l'existence de (16) résulte l'invariance de \(I\) par rapport aux transformations infinitésimales \(\D u,\D x\) de \(\G_{\infty\rho}\) pour \(p(x)\) arbitraire, il en résulte en particulier déjà l'invariance par rapport aux transformations infinitésimales de \(\G_\sigma\) qui naissent de celles-ci par spécialisation, et donc par rapport à \(\G_\sigma\). Les relations de divergence
% unit: noether.p13.eq.048
\[
\sum\psi_i\bd u_i^{(\lambda)}=\Div B^{(\lambda)}
\]
doivent donc être des conséquences des dépendances (16), que l'on peut aussi écrire
% unit: noether.p13.eq.049
\[
\sum\psi_i a_i^{(\lambda)}=\Div\chi^{(\lambda)},
\]
où les \(\chi^{(\lambda)}\) sont des combinaisons linéaires des expressions de Lagrange et de leurs dérivées. Comme les \(\psi\) interviennent linéairement aussi bien dans (13) que dans (16), les relations de divergence doivent donc, en particulier, être des combinaisons linéaires des dépendances (16). Ainsi
% unit: noether.p13.eq.050
\[
\Div B^{(\lambda)}=\Div\left(\sum \alpha\cdot\chi^{(\varkappa)}\right),
\]
et les \(B^{(\lambda)}\) eux-mêmes se composent donc linéairement des \(\chi\), c'est-à-dire des expressions de Lagrange et de leurs dérivées, et de fonctions dont la divergence s'annule identiquement, comme par exemple les \(B-\Gamma\) qui apparaissent à la fin du \S{}2, pour lesquels \(\Div(B-\Gamma)=0\), et où la divergence possède en même temps une propriété d'invariant. J'appellerai \og impropres \fg{} les relations de divergence dans lesquelles les \(B^{(\lambda)}\) peuvent se composer de la manière indiquée à partir des expressions de Lagrange et de leurs dérivées; toutes les autres seront appelées \og propres \fg{}.

% unit: noether.p13.par.071
Inversement, si les relations de divergence sont des combinaisons linéaires des dépendances (16), donc \og impropres \fg{}, alors l'invariance par rapport à \(\G_\sigma\) résulte de celle par rapport à \(\G_{\infty\rho}\); \(\G_\sigma\) devient un sous-groupe de \(\G_{\infty\rho}\). Les relations de divergence correspondant à un groupe fini \(\G_\sigma\) deviennent donc impropres si et seulement si \(\G_\sigma\) est un sous-groupe d'un groupe infini par rapport auquel \(I\) est invariant.

% unit: noether.p13.par.072
Par spécialisation des groupes, l'assertion originale de Hilbert s'en déduit. Par \og groupe des translations \fg{}, on entend le groupe fini
% unit: noether.p13.eq.051
\[
y_i=x_i+\eps_i,\qquad v_i(y)=u_i(x);
\]
donc
% unit: noether.p13.eq.052
\[
\D x_i=\eps_i,\qquad \D u_i=0,\qquad
\bd u_i=-\sum_\lambda\frac{\p u_i}{\p x_\lambda}\eps_\lambda.
\]
L'invariance par rapport au groupe des translations signifie, comme on sait, que dans
% unit: noether.p13.eq.053
\[
I=\int\!\cdots\!\int f\!\left(x,u,\frac{\p u}{\p x},\ldots\right)\dd x
\]
les \(x\) n'apparaissent pas explicitement dans \(f\). Les \(n\) relations de divergence correspondantes
% unit: noether.p13.eq.054
\[
\sum\psi_i\frac{\p u_i}{\p x_\lambda}=\Div B^{(\lambda)}
\qquad(\lambda=1,2,\ldots,n)
\]
seront appelées \og relations d'énergie \fg{}, parce que les \og lois de conservation \fg{} \(\Div B^{(\lambda)}=0\) qui correspondent au problème variationnel répondent aux \og théorèmes d'énergie \fg{}, et les \(B^{(\lambda)}\) aux \og composantes d'énergie \fg{}. On obtient donc: si \(I\) admet le groupe des translations, les relations d'énergie deviennent impropres si et seulement si \(I\) est invariant par rapport à un groupe infini contenant le groupe des translations comme sous-groupe.\footnote{Les théorèmes d'énergie de la mécanique classique, de même que ceux de l'ancienne \og théorie de la relativité \fg{} (où \(\sum \dd x^2\) se transforme en lui-même), sont \og propres \fg{}, car aucun groupe infini n'intervient ici.}

% unit: noether.p13.par.073
Un exemple de tels groupes infinis est donné par le groupe de toutes les transformations des \(x\) et des transformations induites des \(u(x)\), dans lesquelles n'apparaissent que des dérivées des fonctions arbitraires \(p(x)\). Le groupe des translations naît par la spécialisation \(p^{(i)}(x)=\eps_i\); il faut toutefois laisser indécise la question de savoir si, par ce moyen --- et par les groupes qui naissent lorsqu'on modifie \(I\) par une intégrale de bord --- on a déjà donné les groupes les plus généraux de ce type. Des transformations induites du genre indiqué apparaissent par exemple lorsqu'on soumet les \(u\) aux transformations de coefficients d'une \og forme différentielle totale \fg{}, c'est-à-dire d'une forme
% unit: noether.p13.eq.055
\[
\sum a\,\dd^\lambda x_i+\sum b\,\dd^{\lambda-1}x_i\,\dd x_\varkappa+\cdots,
\]
qui contient, en plus des \(\dd x\), des différentielles d'ordre supérieur. Des transformations induites plus spéciales, dans lesquelles les \(p(x)\) n'apparaissent qu'en première dérivée, sont données par les transformations de coefficients des formes différentielles ordinaires \(\sum c\,\dd x_{i_1}\cdots \dd x_{i_\lambda}\); ce sont celles qu'on a ordinairement considérées.

% unit: noether.p13.par.074
Un autre groupe du type indiqué --- qui, à cause de la présence du terme logarithmique, ne peut pas être une transformation de coefficients --- est par exemple le suivant:
% unit: noether.p13.eq.056
\[
y=x+p(x),\qquad
v_i=u_i+\lg(1+p'(x))=u_i+\lg\frac{\dd y}{\dd x};
\]
% unit: noether.p13.eq.057
\[
\D x=p(x),\qquad
\D u_i=p'(x)\srcfnmark{1)},\qquad
\bd u_i=p'(x)-u_i'p(x).
\]
\srcfntext{1)}{À partir de ces transformations infinitésimales, on calcule les transformations finies en remontant selon la méthode indiquée à la fin du \S{}4.}
Les dépendances (16) deviennent ici:
% unit: noether.p13.eq.058
\[
\sum_i\left(\psi_i u_i'+\frac{\dd\psi_i}{\dd x}\right)=0,
\]
et les relations d'énergie impropres deviennent:
% unit: noether.p13.eq.059
\[
\sum\left(\psi_i u_i'+
\frac{\dd(\psi_i+\mathrm{const.})}{\dd x}\right)=0.
\]
Une intégrale invariante la plus simple du groupe est:
% unit: noether.p13.eq.060
\[
I=\int \frac{e^{-2u_1}}{u_1'-u_2'}\,\dd x.
\]
L'intégrale \(I\) la plus générale se détermine par intégration de l'équation différentielle de Lie (11)
% unit: noether.p13.eq.061
\[
\bd f+\frac{\dd}{\dd x}(f\cdot \D x)=0,
\]
qui, après substitution des valeurs de \(\D x\) et de \(\bd u\), devient, dès que l'on suppose \(f\) dépendant seulement des dérivées premières des \(u\),
% unit: noether.p13.eq.062
\[
\frac{\p f}{\p x}p(x)
+\left\{\sum\frac{\p f}{\p u_i}
-\frac{\p f}{\p u_i'}u_i'+f\right\}p'(x)
+\left\{\sum\frac{\p f}{\p u_i'}\right\}p''(x)=0
\]
(identiquement en \(p(x),p'(x),p''(x)\)). Déjà pour deux fonctions \(u(x)\), ce système d'équations possède des solutions qui contiennent effectivement les dérivées, à savoir
% unit: noether.p13.eq.063
\[
f=(u_1'-u_2')\,
\Phi\!\left(u_1-u_2,\frac{e^{-u_1}}{u_1'-u_2'}\right),
\]
où \(\Phi\) désigne une fonction arbitraire des arguments indiqués.

% unit: noether.p13.par.075
Hilbert énonce son assertion en disant que l'échec des théorèmes d'énergie propres est un caractère distinctif de la \og théorie générale de la relativité \fg{}. Pour que cette assertion vaille littéralement, il faut donc comprendre la désignation \og relativité générale \fg{} plus largement que d'ordinaire, et l'étendre aussi aux groupes précédents dépendant de \(n\) fonctions arbitraires.\footnote{Ainsi se confirme de nouveau la justesse d'une remarque de Klein, selon laquelle la désignation \og relativité \fg{} usuelle en physique devrait être remplacée par \og invariance relativement à un groupe \fg{}. (\emph{Über die geometrischen Grundlagen der Lorentzgruppe}. \emph{Jhrber. d. d. Math. Vereinig.} vol. 19, p. 287, 1910; réimprimé dans la \emph{Physikalische Zeitschrift}.)}
